% ===========================================================================
\section{Introduction}\label{sec:intro}
% ===========================================================================

A \emph{Butson Hadamard matrix} $\BH(n,h)$ \cite{Butson1962} is an
$n\times n$ matrix $H$ with entries in the group $\mu_h$ of complex
$h$-th roots of unity and $HH^*=nI_n$. For a finite abelian group $G$, a
\emph{$G$-invariant} $\BH(G,h)$ has entries $f(g'-g)$ for a function
$f:G\to\mu_h$. For $G=C_n$ these are the circulant $\BH(n,h)$, whose
first rows are the \emph{perfect sequences} of length $n$ over $\mu_h$;
equivalent objects are generalized bent functions $G\to\Z_h$ and, for
prime $h$, semi-regular relative difference sets. Since Turyn
\cite{Turyn1965}, nonexistence proofs work with the ideal factorization
of the character values of $X=\sum_gf(g)g\in\Z[\zeta_h][G]$, which have
absolute value $\sqrt{|G|}$: self-conjugacy, Ma's lemma
\cite{Ma1985,Pott1995}, field descent \cite{Schmidt1999}, class groups.

\subsection{Two conjectures}\label{ssec:conj}
Let $r$ be the square-free part of $n$. The following conjecture is
attributed to Mow \cite{Mow1995} in K.-U.~Schmidt's survey
\cite[Conj.~4.1.1]{Schmidt2016survey} (we could not consult the text of
\cite{Mow1995}, whose abstract describes a unified construction):
\emph{a perfect sequence of length $n>1$ over $\mu_H$ exists if and only
if $H$ is divisible by $2\sqrt{rn}$ for $n\equiv2\pmod4$ and by
$\sqrt{rn}$ otherwise.} Do Duc \cite[Conj.~3.1]{DoDuc2019}\footnote{We
number the results of \cite{DoDuc2019} as in arXiv:1903.07334v2.}
conjectured:
\emph{a $\BH(\Z_n,h)$ exists if and only if (i)
$\nu_p(h)\ge\lceil\nu_p(n)/2\rceil$ for every prime $p\mid n$, and (ii)
$\nu_2(h)\ge2$ if $n\equiv2\pmod4$.} Since
$\nu_p(\sqrt{rn})=\lceil\nu_p(n)/2\rceil$, the two conjectures are the
same statement. Sufficiency follows from the constructions of Mow and of
Duc and Schmidt \cite[Cor.~2.5]{DucSchmidt2019}, \cite[Result~4.1]{DoDuc2019}. Do Duc writes that ``we are
far from proving Conjecture~3.1'' and lists $2687$ pairs $(n,h)$,
$n,h\le100$, as open \cite[Remark~3.13]{DoDuc2019}.

\subsection{The main theorem}

\begin{theorem}\label{thm:S}
Let $G$ be a finite abelian group and $p$ a prime with $p\,\|\,|G|$. If
$p$ is unramified in $\Q(\zeta_h)$, that is, $p\nmid h$, or $p=2$ and
$4\nmid h$, then there is no $\BH(G,h)$. In particular, if a perfect
sequence of length $n$ over $\mu_h$ exists, then every prime $p$ with
$\nu_p(n)=1$ divides $h$, and $4\mid h$ if $n\equiv2\pmod4$.
\deptag{Lean}
\end{theorem}

The proof (Section~\ref{sec:S}) writes $G=H\times C_p$, evaluates only
the characters of $H$, and keeps $C_p$ as a group-ring variable. Modulo a
prime above $p$ the group ring of $C_p$ becomes the local ring
$\F[T]/(T-1)^p$; this forces all coefficients of the partial evaluation
into one prime above $p$, and Parseval's identity on a single coset of
$C_p$ gives a contradiction.

\subsection{Consequences}
(a) Conditions (i) at every prime $p$ with $\nu_p(n)=1$ and (ii) are
necessary for every abelian group of order $n$ (Corollary~\ref{cor:doduc}).
With the known sufficiency this gives:

\begin{corollary}\label{cor:squarefree}
Let $n$ and $h$ be positive integers. If $n$ is square-free, a
$\BH(\Z_n,h)$ exists if and only if $n\mid h$
and, for even $n$, $4\mid h$. Thus the conjectures of Mow and of Do Duc
hold for square-free $n$; e.g.\ a circulant $\BH(2p,h)$, $p$ an odd
prime, exists if and only if $4p\mid h$. \deptag{Lean-derived}
\end{corollary}

\noindent Theorem~\ref{thm:S} excludes $2323$ of the $2687$ pairs listed
in \cite[Remark~3.13]{DoDuc2019}; $2063$ of them are not excluded by any
result known to us, subject to \cite{LeungChueZhao2025}. (b) For $q\equiv2\pmod4$ there is no generalized bent
function $\Z_q\to\Z_q$, the case $m=1$ of the problem of Kumar, Scholtz
and Welch \cite{KumarScholtzWelch1985}, treated so far family by family,
most recently for $q=2\cdot3^a7^b$ \cite{YingDeng2026}
(Corollary~\ref{cor:ksw}). (c) If $D$ is a semi-regular RDS in an
abelian group $G$ relative to $N$ and $p\,\|\,|G/N|$, then $N$ is a
$p$-group, without any splitting assumption (Corollary~\ref{cor:R}).
This contains all nonexistence results of
\cite[Sect.~6]{LeungSchmidtZhang2023} and decides eight entries marked
``?'' in its tables. (d) The eleven orders $260$, $340$, $442$, $468$,
$520$, $580$, $680$, $754$, $820$, $884$, $890$, the only orders $\le1000$
for which a circulant complex Hadamard matrix had not been excluded
according to Arasu's survey \cite{Arasu2011}, are excluded
(Corollary~\ref{cor:arasu}). (e) For perfect quaternary sequences of
length $N\equiv0\pmod4$, $N\le5000$, only $3528$ and $4900$ remain
under the criteria we implemented (Appendix~\ref{app:quaternary}).

(f) \emph{The ramified regime} (Section~\ref{sec:ramified}). For the
family $\BH(pq^e,p)$, $p\ne q$ primes, Theorem~\ref{thm:S} applies only
if $e=1$ or $p=2$. Theorem~\ref{thm:C} shows that there is no circulant
$\BH(pq^e,p)$ for all $e\ge1$. Earlier results cover the case that $q$ is
self-conjugate modulo $p$ \cite[Thm.~3.3]{DoDuc2019} and further
subfamilies; of the $334$ lengths $pq^e\le2000$ with $p$ odd and $e\ge2$,
$34$ are not covered by the earlier criteria we implemented, among them
$56$ and $99$ from Yayla's table \cite{Yayla2016}. The tool is a
comparison lemma at the ramified prime $p$ (Lemma~\ref{lem:A}), proved by
expanding in the uniformizer $\zeta_p-1$: the digits are polynomials in
the exponent of bounded degree, and comparing degrees with the $p$ points
of $\F_p$ gives the claim.

The proofs use no result of the (not peer-reviewed) OpenAI preprint on
the circulant Hadamard conjecture \cite{OpenAI2026}, in whose Lean
library our formalization is placed.

\paragraph{Conventions}
$C_N$ is the cyclic group of order $N$, $\zeta_t=e^{2\pi i/t}$,
$\nu_p=v_p$ the $p$-adic valuation and $v_\fP$ the valuation at a prime
ideal $\fP$. Tags: \deptag{Lean}: the statement is an instance of a
formally verified Lean theorem (Section~\ref{sec:lean});
\deptag{Lean-derived}: it follows from a formally verified statement by a
short argument that is not formalized; \deptag{Comp}: it rests on a
finite computation (Section~\ref{ssec:comp}).
